\section{An Audit Protocol}
\label{sec:audit}

The diagnostics can be assembled into a protocol for deciding whether a
supracompetitive outcome produced by learning algorithms is collusion. It
needs what a firm testing its own algorithm, or a supervisor with access to
a testing environment, typically has: the ability to retrain the algorithm
and to replay a market from a saved state. Investment firms in the European
Union must already test trading algorithms in such environments before
deployment \citep{rts6}, so the protocol adds questions, not infrastructure.

\begin{description}
\item[Step 1: Benchmark.] Compute the competitive and joint-profit
  benchmarks of the environment and the collusion index $\Delta$. Only an
  outcome with $\Delta$ clearly above zero needs the remaining steps.
\item[Step 2: Feasibility.] Ask whether punishment could sustain the outcome
  at all (Propositions~\ref{prop:trigger} and~\ref{prop:anypun}). An outcome
  above what any punishment can sustain is not equilibrium collusion,
  whatever else is true; one above what Nash reversion can sustain requires
  harsher punishments, which Step~3 would see.
\item[Step 3: Deviation test.] Force a one-period best-response deviation in a
  replayed market and compare it with an undeviated copy under common random
  numbers. Collusion predicts that rivals respond and that deviating does not
  pay over the following periods.
\item[Step 4: Placebos.] Retrain (a) myopically ($\gamma=0$), (b) without
  memory and (c) with an uninformative memory of the same size. None of
  these can sustain punishment, and by Proposition~\ref{prop:myopic} (a)
  plays Nash if its estimates are unbiased. Collusion predicts that the
  outcome, and any reaction to deviations or shocks, disappears.
\item[Step 5: Learning rule.] Retrain with a different step size and with
  counterfactual updates. An equilibrium is a property of the game and should
  not scale with the step size; a learning bias should.
\end{description}

\noindent The verdict is \emph{collusion} if rivals respond to the deviation
and deviating does not pay (Step 3), and the learners that cannot punish
(Step 4) show no such response; \emph{learning bias} if deviating pays and
the outcome is reproduced by a placebo or moves with the learning rule
(Steps 4 and 5); and inconclusive otherwise. The verdict rests on the
response to deviations, not on the level of prices or trading: Steps 4 and 5
show how easily that level is produced without punishment. Each step can
fail, and a genuine punishment scheme passes all of them.

Table~\ref{tab:audit} applies the protocol to every market in this paper. It
returns \emph{collusion} in exactly one case: the \citet{calvano2020}
baseline. Every case among informed traders, and the memoryless pricing
duopoly, returns \emph{learning bias}, although several of them would pass a
test based on the price or trading level alone, and one (the myopic
shared-table learner at $\xi=500$) would pass a test based on the reaction to
price shocks.

\begin{table}[t]
  \centering
  \caption{The audit protocol applied to each market. $\Delta$: profit-based
  (Bertrand, dealers) or intensity-based (Kyle). Deviation: does the rival
  respond and does deviating lose? Placebo: does a learner that cannot punish
  reproduce the outcome? Rule: does the outcome change with the learning
  rule (counterfactual updates or step size)?}
  \label{tab:audit}
  \resizebox{\linewidth}{!}{\begin{tabular}{lllccccl}
\toprule
Market & Learners & $\Delta$ & Rival responds & Deviating & Placebo $\Delta$ (no punishment possible) & Counterfactual $\Delta$ & Verdict \\
\midrule
Logit Bertrand & memory one & $\BertBase$ & yes ($\BertBaseRival$) & loses ($\BertBaseGain$) & myopic $\BertMyopic$, no memory $\BertNone$ & $\BertCf$ & collusion \\
Logit Bertrand & no memory & $\BertNone$ & no & pays ($+\BertNoneGain$) & (is a placebo) & -- & learning bias \\
Kyle, $\xi=0$ & residual memory & $\ResidDeltaInt$ & no & pays ($+\ResidGain$) & myopic $\MyopicResidDeltaInt$, random memory $\RandomMemDeltaInt$ & $\CfResidDeltaInt$ & learning bias \\
Kyle, $\xi=0$ & perfect monitoring & $\OrdersDeltaInt$ & no & pays ($+\OrdersGain$) & as above & -- & learning bias \\
Kyle, $\xi=500$ & separate tables & $\DouStratDeltaInt$ & no & pays & myopic $\DouMyopicDeltaInt$ & -- & learning bias \\
Kyle, $\xi=500$ & shared table & $\SharedDouStratDeltaInt$ & no & pays & myopic $\SharedDouMyopicDeltaInt$ (reacts to shocks) & -- & learning bias \\
Kyle, $\xi=500$ & \citeauthor{dou2025}'s protocol & $\FaDeltaProfit$ (profit) & borderline ($\FaRival\%$) & pays & myopic $\FaMyopicDeltaProfit$ (same shock response); no lagged price $\FaNoPriceDeltaProfit$ & -- & learning bias \\
Dealers & memory one & $\QuoteBaseProfit$ & \QuoteBaseRivalWord & \QuoteBaseGainWord & myopic $\QuoteMyopicProfit$, no memory $\QuoteNoneProfit$ & $\QuoteCfProfit$ & \QuoteVerdict \\
\bottomrule
\end{tabular}
}
\end{table}
